\section{Método}\label{sec:met}

Para la realización de este estudio, la metodología se ha estructurado de la siguiente manera:

\subsection{Formulación Matemática del Problema}\label{subsec:formulacion}

Sea $G = (V, E)$ un grafo no dirigido con conjunto de vértices $V = \{1, 2, \ldots, n\}$ y conjunto de aristas $E$. Se define $\bar{E} = \{(i,j) : i,j \in V,\ i \neq j,\ (i,j) \notin E\}$ como el conjunto de aristas complementarias, es decir, los pares de vértices que \textbf{no} están conectados en $G$.

\subsubsection{Formulación como Programa Lineal Entero (ILP)}

El MCP se formula como un problema de Programación Lineal Entera Binaria (Binary Integer Linear Program, BILP). Se introduce una variable de decisión binaria $x_i$ para cada vértice $i \in V$, donde $x_i = 1$ si el vértice $i$ pertenece al clique, y $x_i = 0$ en caso contrario \cite{pardalos1992_maxclique, nemhauser1988}.

\vspace{0.3cm}
\noindent\textbf{Función objetivo:}
\begin{equation}\label{eq:obj}
    \max \quad z = \sum_{i \in V} x_i
\end{equation}

\noindent\textbf{Restricciones de adyacencia:}
\begin{equation}\label{eq:adj}
    x_i + x_j \leq 1, \quad \forall\ (i,j) \in \bar{E}
\end{equation}

\noindent\textbf{Restricciones de integralidad:}
\begin{equation}\label{eq:bin}
    x_i \in \{0, 1\}, \quad \forall\ i \in V
\end{equation}

La función objetivo \eqref{eq:obj} maximiza el número de vértices seleccionados. Las restricciones \eqref{eq:adj} garantizan que si dos vértices $i$ y $j$ \textbf{no} son adyacentes (es decir, $(i,j) \in \bar{E}$), entonces ambos no pueden pertenecer simultáneamente al clique. Las restricciones \eqref{eq:bin} aseguran que las variables sean binarias. Toda solución factible de este modelo corresponde a un clique en $G$, y la solución óptima es el clique de cardinalidad máxima \cite{pardalos1992_maxclique, bomze1999}.

\subsubsection{Relajación Lineal Continua (LP)}

Para obtener cotas superiores del valor óptimo entero, se resuelve la relajación lineal del problema \eqref{eq:obj}--\eqref{eq:bin}, reemplazando la restricción de integralidad por su envolvente continua \cite{nemhauser1988, schrijver1986}:

\begin{equation}\label{eq:relax}
    0 \leq x_i \leq 1, \quad \forall\ i \in V
\end{equation}

Sea $z^*_{LP}$ el valor óptimo de la relajación y $z^*_{ILP}$ el del problema entero original. Dado que la región factible del LP contiene a la del ILP, se cumple la relación fundamental:
\begin{equation}\label{eq:cota}
    z^*_{ILP} \leq z^*_{LP}
\end{equation}

Esta propiedad es la base teórica del método de Ramificación y Acotamiento: si la cota superior $z^*_{LP}$ de un subproblema no supera la mejor solución entera encontrada hasta el momento, dicho subproblema puede descartarse sin pérdida de optimalidad (\textit{poda por cota}) \cite{nemhauser1988, wolsey1998}.

\subsubsection{Formulación Cuadrática Continua de Motzkin-Straus}

Una formulación alternativa del MCP proviene del teorema de Motzkin-Straus \cite{motzkin1965}, que establece una equivalencia exacta entre el problema combinatorio discreto y un problema de programación cuadrática continua. Sea $A$ la matriz de adyacencia del grafo $G$ y $\Delta$ el simplex estándar $\Delta = \{x \in \mathbb{R}^n : \sum_{i=1}^{n} x_i = 1,\ x_i \geq 0\}$. El teorema establece que:

\begin{equation}\label{eq:motzkin}
    \frac{1}{2}\left(1 - \frac{1}{\omega(G)}\right) = \max_{x \in \Delta} \quad \frac{1}{2} x^T A x
\end{equation}

\noindent donde $\omega(G)$ denota el tamaño del clique máximo de $G$. Esta formulación transforma el problema discreto en uno de optimización continua no lineal, lo cual permite emplear técnicas de programación cuadrática para su resolución \cite{pardalos1992_maxclique, bomze1999}.

\subsection{Procedimiento de Resolución}

\begin{enumerate}
    \item \textbf{Herramientas de desarrollo:} Para los experimentos, el equipo computacional ejecutará algoritmos implementados en el lenguaje de programación Python. Python ofrece una robusta capacidad para la formulación rápida de modelos de optimización y el manejo eficiente de matrices de adyacencia. 
    
    \item \textbf{Materiales de prueba:} Los materiales consistirán en grafos estándar extraídos de bancos de pruebas (\textit{benchmarks}) ampliamente reconocidos en la literatura para evaluar la búsqueda exacta de cliques, tales como los grafos del \textit{2nd DIMACS Challenge} o las instancias de la librería BHOSLIB \cite{dimacs1993, bhoslib, pardalos1992}.
    
    \item \textbf{Procedimiento algorítmico central:} El procedimiento se basará en el método exacto de Ramificación y Acotamiento (\textit{Branch \& Bound}) \cite{carraghan1990, pardalos1992}. 
    
    \item \textbf{Cálculo de la relajación:} El algoritmo comenzará en cada paso resolviendo la relajación lineal o continua del problema original de programación entera \cite{nemhauser1988}. 
    
    \item \textbf{Proceso de ramificación:} Si la solución óptima del problema relajado proporciona un vértice con valores fraccionarios, el procedimiento continuará ramificando sobre una variable discreta $x_k$ que no posea un valor entero \cite{carraghan1990, nemhauser1988}. Dado que el problema se modela con variables binarias ($x_j \in \{0,1\}$), en cada ramificación se fijará el valor de la variable creando dos subproblemas mutuamente excluyentes: $x_k = 0$ y $x_k = 1$
\end{enumerate}

\begin{algorithm}
\caption{Ramificación y Acotamiento (Branch \& Bound) para el MCP}
\label{alg:branch_bound_mcp}
\begin{algorithmic}[1]
\Require Grafo $G=(V,E)$, Formulación binaria del problema (Maximización)
\Ensure Tamaño del clique máximo $\check{z}$ y el vector de solución $x^*$
\State $L \gets \{ P_0 \}$ \Comment{Inicializar la lista de nodos con el problema original}
\State $\check{z} \gets 0$ \Comment{Mejor cota inferior encontrada (tamaño máximo del clique)}
\State $x^* \gets \emptyset$ \Comment{Vector solución óptimo}
\While{$L \neq \emptyset$}
    \State Seleccionar y remover un subproblema $P$ de la lista $L$
    \State Resolver la relajación lineal de $P$ para obtener la solución $x'$ y su valor $z'$
    \If{$P$ es infactible}
        \State \textbf{continue} \Comment{Poda por infactibilidad}
    \EndIf
    \If{$z' \le \check{z}$}
        \State \textbf{continue} \Comment{Poda por cota: el nodo no puede mejorar la mejor solución}
    \EndIf
    \If{$x'$ es completamente entera}
        \State $\check{z} \gets z'$ \Comment{Actualizar la mejor solución encontrada}
        \State $x^* \gets x'$
    \Else
        \State Seleccionar una variable fraccionaria $x_k$ del vector $x'$
        \State Crear subproblema $P_0$ añadiendo la restricción $x_k = 0$ a $P$
        \State Crear subproblema $P_1$ añadiendo la restricción $x_k = 1$ a $P$
        \State $L \gets L \cup \{ P_0, P_1 \}$ \Comment{Ramificar el espacio de búsqueda}
    \EndIf
\EndWhile
\State \Return $\check{z}, x^*$
\end{algorithmic}
\end{algorithm}